\chapter{Diffusion Models: Mathematical Foundations}

\section{Forward and Reverse Processes}

Diffusion models are built upon a pair of stochastic processes: a forward (noising) process and a reverse (denoising) process.

\medskip

\noindent
\textbf{Forward process.}  
We define a Markov chain $(x_0, x_1, \dots, x_T)$ such that:
\[
q(x_t \mid x_{t-1}) = \mathcal{N}(x_t; \sqrt{1 - \beta_t}\, x_{t-1}, \beta_t I),
\]
where $\{\beta_t\}_{t=1}^T$ is a variance schedule.

\medskip

\noindent
\textbf{Closed-form expression.}  
By composing Gaussian transitions, we obtain:
\[
q(x_t \mid x_0)
=
\mathcal{N}(x_t; \sqrt{\bar{\alpha}_t}\, x_0, (1 - \bar{\alpha}_t) I),
\]
where $\alpha_t = 1 - \beta_t$ and $\bar{\alpha}_t = \prod_{s=1}^t \alpha_s$.

\medskip

\noindent
\textbf{Interpretation.}
\begin{itemize}
    \item As $t$ increases, $x_t$ becomes increasingly noisy
    \item For large $T$, $x_T \approx \mathcal{N}(0, I)$
\end{itemize}

\medskip

\noindent
\textbf{Reverse process.}  
We seek to model:
\[
p_\theta(x_{t-1} \mid x_t),
\]
which is also Gaussian under suitable assumptions:
\[
p_\theta(x_{t-1} \mid x_t)
=
\mathcal{N}(x_{t-1}; \mu_\theta(x_t, t), \Sigma_\theta(x_t, t)).
\]

\medskip

\noindent
\textbf{Learning objective.}  
The parameters $\theta$ are trained to approximate the true reverse transitions.

\medskip

\noindent
\textbf{Insight.}  
The forward process is fixed and simple, while the reverse process is learned.

\section{Score Matching}

An alternative formulation of diffusion models is based on score matching.

\medskip

\noindent
\textbf{Score function.}
\[
s_t(x) = \nabla_x \log p_t(x),
\]
where $p_t(x)$ is the distribution of $x_t$.

\medskip

\noindent
\textbf{Objective.}  
Instead of modeling $p_t(x)$ directly, we learn $s_t(x)$.

\medskip

\noindent
\textbf{Denoising score matching.}  
For Gaussian perturbations, the score satisfies:
\[
s_t(x_t) = -\frac{1}{1 - \bar{\alpha}_t} (x_t - \sqrt{\bar{\alpha}_t} x_0).
\]

\medskip

\noindent
\textbf{Learning formulation.}  
We train a neural network $s_\theta(x_t, t)$ to approximate the score:
\[
\mathbb{E}_{t, x_0, \varepsilon}
\left[
\| s_\theta(x_t, t) - \nabla_{x_t} \log q(x_t \mid x_0) \|^2
\right].
\]

\medskip

\noindent
\textbf{Equivalence.}  
This objective is equivalent to predicting the noise $\varepsilon$.

\medskip

\noindent
\textbf{Insight.}  
Score matching provides a principled way to learn gradients of log-densities.

\section{Langevin Dynamics}

Sampling from a distribution can be achieved using Langevin dynamics.

\medskip

\noindent
\textbf{Langevin update.}
\[
x_{k+1} = x_k + \eta \, \nabla_x \log p(x_k) + \sqrt{2\eta} \, \xi_k,
\quad \xi_k \sim \mathcal{N}(0, I).
\]

\medskip

\noindent
\textbf{Interpretation.}
\begin{itemize}
    \item Gradient term drives samples toward high-density regions
    \item Noise term ensures exploration
\end{itemize}

\medskip

\noindent
\textbf{Connection to diffusion models.}
\begin{itemize}
    \item Learned score approximates $\nabla_x \log p(x)$
    \item Sampling uses score-based updates
\end{itemize}

\medskip

\noindent
\textbf{Continuous-time limit.}  
Langevin dynamics corresponds to a stochastic differential equation:
\[
dx_t = \nabla_x \log p(x_t)\, dt + \sqrt{2} \, dW_t,
\]
where $W_t$ is Brownian motion.

\medskip

\noindent
\textbf{Insight.}  
Diffusion sampling can be viewed as approximate Langevin dynamics with learned scores.

\section{Sampling Algorithms}

Several sampling procedures arise from the mathematical formulation of diffusion models.

\medskip

\noindent
\textbf{Ancestral sampling.}
\begin{itemize}
    \item Start from $x_T \sim \mathcal{N}(0, I)$
    \item Iteratively sample:
    \[
    x_{t-1} \sim p_\theta(x_{t-1} \mid x_t)
    \]
\end{itemize}

\medskip

\noindent
\textbf{Deterministic variants.}
\begin{itemize}
    \item Use deterministic updates instead of sampling
    \item Reduce variance and speed up generation
\end{itemize}

\medskip

\noindent
\textbf{Score-based sampling.}
\begin{itemize}
    \item Use discretized Langevin dynamics
    \item Update using learned score function
\end{itemize}

\medskip

\noindent
\textbf{Tradeoffs.}
\begin{itemize}
    \item More steps $\rightarrow$ higher quality samples
    \item Fewer steps $\rightarrow$ faster generation
\end{itemize}

\medskip

\noindent
\textbf{Numerical considerations.}
\begin{itemize}
    \item Choice of step size
    \item Stability of updates
\end{itemize}

\medskip

\noindent
\textbf{Insight.}  
Sampling is an iterative refinement process guided by learned gradients.

\section{A Unifying Perspective}

Diffusion models can be understood through multiple mathematical frameworks:

\begin{itemize}
    \item \textbf{Markov processes:} forward and reverse chains
    \item \textbf{Score matching:} learning gradients of log-density
    \item \textbf{Stochastic dynamics:} Langevin equations
\end{itemize}

\medskip

\noindent
\textbf{Core components.}
\begin{itemize}
    \item Forward noise process $q(x_t \mid x_{t-1})$
    \item Reverse process $p_\theta(x_{t-1} \mid x_t)$
    \item Score function $s_\theta(x_t, t)$
\end{itemize}

\medskip

\noindent
\textbf{Key insight.}  
Diffusion models unify probabilistic modeling, stochastic processes, and gradient-based sampling.

\section{Outlook}

This chapter developed the mathematical foundations of diffusion models:
\begin{itemize}
    \item forward and reverse stochastic processes,
    \item score matching formulations,
    \item Langevin dynamics,
    \item sampling algorithms.
\end{itemize}

\medskip

In the next chapter, we study modern refinements and practical implementations of diffusion models.

\medskip

\noindent
\textbf{Guiding principle.}  
Complex distributions can be sampled by following learned gradients within a stochastic dynamical system.